\SourcePage{401}{404}
\section{Symmetry of molecular terms}

In \S~78 we have already examined certain symmetry properties of the terms of a diatomic molecule. These properties described the behavior of the wavefunctions under transformations that leave the nuclear coordinates unaffected. Thus, molecular symmetry under reflection in a plane containing its axis accounts for the distinction between the $\Sigma^+$ and $\Sigma^-$ terms; symmetry under reversal of the signs of all electron coordinates (for a molecule made up of identical atoms)\footnote{The coordinate origin is assumed to lie on the molecular axis, midway between the two nuclei.} leads to the division of the terms into even and odd ones.
